\scanpage{15}
\mspage{PU 74}

Par sa ``texture sensorielle'', l'image\deleted{\unsure{ sym}} peut consister en des sons, \handwritten{ou} des odeurs, \handwritten{ou} des goûts, ou en des sensations tactiles de toute nature et perçues en n'importe quelle partie du corps, tout comme en des perceptions visuelles. Dans la plupart des rêves, plusieurs de nos sens sont mis à contribution en même temps. Mais cette ``texture\deleted{\unsure{''}} de fond'' que fournit la sensation n'épuise et ne ``dit'' pas plus l'``image'', qu'un rythme ne dit une mélodie, ou que le contour d'une fleur ne dit le jeu délicat de la couleur et son mouvement dans la brise, ni ce parfum qui n'est qu'à elle, ni l'ivresse de l'abeille qui vient la butiner. L'émotion dont une image est chargée, et \handwritten{\inserted{aussi}} le nuage plus ou moins net ou plus ou moins diffus des associations qui \handwritten{l'}entourent\handwritten{,}\deleted{l'image,} font partie d\handwritten{u} \deleted{son} sens \handwritten{\inserted{de l'image}} et de son message, aussi intimement que le parfum fait partie d'une fleur.

La langue-images reste ainsi\handwritten{,} à tout moment\handwritten{,} en contact immédiat avec la perception des \mbox{s\,e\,n\,s} (\textsuperscript{*}), avec le \mbox{c\,o\,r\,p\,s}\handwritten{\inserted{(**)}}; et aussi avec l'\uline{émotion}, fille du corps et des sens, fidèle messagère de ce qui est perçu.

\deleted{Ainsi} \mbox{\handwritten{P}\,e\,r\,c\,e\,p\,t\,i\,o\,n}, \mbox{é\,m\,o\,t\,i\,o\,n}, et \mbox{e\,x\,p\,r\,e\,s\,s\,i\,o\,n} d' \mbox{u\,n\quad s\,e\,n\,s} (ou d'une `` \mbox{p\,e\,n\,s\,é\,e} '') sont \deleted{\unsure{tous}} ici une seule et même chose.

% slanted
\margin{\handwritten{\unsure{espacer}}}

\deleted{Je me \inserted{suis} laissé\deleted{\unsure{nt}} entraîner (\deleted{\unsure{par le nom dans}}\handwritten{\inserted{tombant dans le}} piège du nom même ``langue des images'' !) à parler de ``l'image'' comme si c'était une entité séparée, un ``élément'' d'expression d'un\deleted{\unsure{x}} ``sens'' ou d'une ``pensée'' - comme le ``mot'' est pierre de construction d'un ``discours''. Pourtant, il n'y a \handwritten{\inserted{ici}} pas plus de ``mots'' s'articulant en ``phrases'', ni même d'``images'' s'assemblant pour former un ``sens'', qu'il n'y en a dans le vol d'une mouette, dans \deleted{le}\inserted{l'incessant} mouvement de l'eau d'un fleuve ou d'un ruisseau, dans la danse d'une libellule. Pour mieux dire encore, à chaque instant, le ``récit'' \mbox{e\,s\,t} ce vol, ce mouvement, cette danse, il est \mbox{v\,i\,e}, vécue au fil des instants par celui qui, en vivant \deleted{\unsure{allégrement|allègrement}} vol, mouvement ou danse et sans même le savoir ni le ``vouloir'', ``fait le \deleted{\unsure{``}}récit''\inserted{(**)}. Très rarement m'est il arrivé que le ``récit'' ou la ``pensée'' \deleted{d'un rêve} exprimés par un rêve tienne en une seule image plus ou moins stati\deleted{\unsure{t}}que, \handwritten{ou} \deleted{même} dans une \handwritten{\inserted{simple}} succession de \deleted{quelques}\handwritten{\inserted{telles}} images. Plutôt, récit et sens sont \deleted{\unsure{alors dits}}\handwritten{\inserted{à chaque fois}} \mbox{m\,i\,s\quad e\,n\quad s\,c\,è\,n\,e} par une sorte de ``psychodrame'', \handwritten{\inserted{par}} une \mbox{p\,a\,r\,a\,b\,o\,l\,e} plus ou moins transparente. \deleted{Ce ``psychodrame'', \unsure{ou}}\handwritten{\inserted{Et avant toute autre chose,}} ce récit-parabole, nous parle \deleted{des forces avant tout} de nous-même: des forces ignorées qui reposent en nous, \deleted{\unsure{ceux}}\inserted{et de leur travail souterrain, et} des conflits, tensions, mascarades, \handwritten{\inserted{\unsure{évènements|événements}}} se jouant dans notre être. \handwritten{D}e ce qui \uline{est}, \handwritten{\inserted{(\deleted{\illegible{1}} que nous évitons de voir\ldots),}} \deleted{comme} de ce qui \uline{fût} (\deleted{\unsure{que}}\handwritten{\inserted{que nous avons}} depuis longtemps oublié\ldots), et d'un ``\uline{possible}'' insoupçonné qui nous attend, (et qu'\handwritten{il} ne tient qu'à nous \handwritten{de réaliser \ldots}).}

% archive sheet number: 14
